% Page 036 — page imprimée 32
% Contenu seul : le préambule et la pagination sont gérés par meyrueis.tex

\begin{center}
{\large\bfseries LA GROTTE DE « LA DOUCE »}
\end{center}

\vspace{1.5em}

\textit{Son entrée, au dessus du Ferradou, sur le vieux chemin montant au rocher du Château, servit
longtemps de bergerie pour quelques brebis. C’est seulement en 1920 que Frank Loupiac et
son jeune cousin, Daniel Couderc\footnote{Daniel Couderc 1908-1997, fils de Léon Couderc et Madeleine de Felice.
Cousin germain de Frank Loupiac}, découvrirent, les premiers, le début du réseau de la
Douce\footnote{La Douce signifie « la Source ». cf André Fages : « La quête de l’eau » 2004 «
Los Adrelhans »} qu’on appelle maintenant grotte de la Porte. Philippe Chambon en est propriétaire et
dispose des clefs.}

\textit{Voici le récit de cette première « exploration » tel qu’elle nous a été contée par Daniel
Couderc et enregistrée au Moulin d’Ayres au cours de l’été 1981.}

\vspace{1em}

« Comme la plupart des enfants j’ai longtemps aimé l’obscurité. Quoi de plus rassurant
lorsqu’on a deux ans de se blottir sous les couvertures tout contre la chaleur maternelle, la
conscience épanouie par de vagues et heureuses réminiscences. Plus tard vers 4 ans, vers 6
ans l’aventure se complique : elle consiste à se frayer soi-même un chemin vers la lumière,
vers le bout du lit étroitement bordé « à la française » : plus tu avances et moins il y a de
place. La tête sert de butoir, elle écarte et soulève les draps pour trouver un peu d’air, il fait
chaud, tu étouffes... mais voilà que ton bras s’est glissé au dehors, il te tire vers la fraîcheur...
Encore un effort et tu débouches en riant tout près de tes parents dans la chambre ensoleillée
par une matinée de printemps. C’est plus tard que pour beaucoup, l’obscurité se peuple d’êtres
malfaisants : fantômes, vampires, araignées ou reptiles venus de légendes racontées par des
adultes irresponsables !!!

Il est possible que ces « histoires » aussi vieilles que l’humanité ne soient elles mêmes que le
reflet d’une crainte ancestrale venue des temps où l’obscurité était en effet pleine de pièges.
Mais c’est là une autre question...
Heureusement, et grâce à des parents éclairés, notre imagination fut préservée de ces craintes
irraisonnées, de même que celle de mon cousin Frank.

Bien qu’il soit mon aîné de 3 ans, nous nous entendions à merveille lorsque je le retrouvais à
chacun de mes séjours à Meyrueis. Il me montrait la rivière ou plutôt les trois rivières qui se
rencontrent à Meyrueis.
Pêcheur émérite il en connaissait tous les « gouffres », et lorsqu’il partait en expédition,
parfois pour toute la journée, il ne revenait jamais bredouille. Mais il aimait aussi les trous, les
grottes, les avens qui abondent dans la région et là, je le suivais avec passion.

En cette année 1920, j’avais 12 ans, lui 15 nous avions décidé de réaliser un projet qui nous
tentait depuis longtemps : nous voulions explorer la grotte de la Douce. La Douce n’est pas
une rivière, même pas un ruisseau, c’est seulement une cascade parfois impétueuse qui se
manifeste après une période de mauvais temps et de pluies abondantes. Elle jaillit du rocher et
tombe presque directement dans la Jonte juste derrière notre maison. Le trou d’où elle sort
parfaitement sec en été est trop étroit pour être remonté au-delà d’une dizaine de mètres mais
son débit lorsqu’elle fonctionne montre qu’elle est le déversoir de tout un réseau de galeries
souterraines s’étendant sous le Causse. Comment atteindre ce réseau ?»
